\makeatletter
\def\input@path{{./}{../}{../../}{../../../}{../../../../}{preamble/}{../preamble/}{../../preamble/}{../../../preamble/}}
\makeatother

\input{preamble_exams.tex}

\title{Grade 11 -- Term 1 \\ Classroom Activity: Sample Statistics}
\author{}
\date{}

\pagestyle{empty}

% ============================================================
% SOLUTION STAGE COMMAND
% ============================================================

\newcommand{\SolutionStage}[2]{%
	\Needspace{6\baselineskip}%
	\subsection{#2}%
}

\begin{document}
	
	% ============================================================
	% PROBLEM PAGES
	% ============================================================
	
	\newgeometry{
		left=0.45in,
		right=0.45in,
		top=0.35in,
		bottom=0.35in
	}
	
	\vspace*{-15mm}
	
	\begin{center}
		{\large\bfseries G11 T1 C7 -- Sample Mean, Variance, and Standard Deviation}
	\end{center}
	
	\vspace{-3mm}
	
	\noindent
	\textbf{Required equations.}
	For a sample $x_1,x_2,\ldots,x_n$, use the sample statistics
	\[
	\boxed{
		n=\text{number of observations},\qquad
		\bar{x}=\frac{\sum_{i=1}^{n}x_i}{n},\qquad
		s^2=\frac{\sum_{i=1}^{n}(x_i-\bar{x})^2}{n-1},\qquad
		s=\sqrt{s^2}
	}
	\]
	\[
	\textrm{Corresponding theoretical (population) values:}
	\qquad
	\boxed{
		\mu=E(X),\qquad
		\sigma^2=Var(X),\qquad
		\sigma=\sqrt{\sigma^2}
	}
	\]
	\[
	\textrm{For a value--frequency table}
	\qquad
	n=\sum f,
	\qquad
	\bar{x}=\frac{\sum fx}{n},
	\qquad
	s^2=\frac{\sum f(x-\bar{x})^2}{n-1}.
	\]
	
	\medskip
	\noindent
	\textbf{Instructions.} For every sample, calculate $n$, $\bar{x}$, $s^2$, and
	$s$. Then compare each statistic with the given theoretical value
	$\mu$, $\sigma^2$, or $\sigma$. Round decimal answers to three decimal places.
	
	\begin{multicols}{2}
		\begin{enumerate}[
			leftmargin=*,
			itemsep=7pt,
			parsep=0pt,
			topsep=4pt
			]
			
			\item
			\textbf{Battery operating time (individual data).}
			The operating times, in hours, of five randomly selected batteries are
			\[
			8,\quad 9,\quad 10,\quad 11,\quad 12.
			\]
			The manufacturer reports $\mu=10$ hours and $\sigma^2=4$ hours$^2$.
			
			\item
			\textbf{Local delivery time (individual data).}
			Six randomly selected delivery times, in minutes, are
			\[
			18,\quad 20,\quad 21,\quad 19,\quad 22,\quad 20.
			\]
			Historical records give $\mu=20$ minutes and $\sigma=1.5$ minutes.
			
			\item
			\textbf{Reusable cups (value--frequency table).}
			The number of reusable cups brought by each student in a random sample is
			summarized below. The school-wide model has $\mu=2.5$ cups and
			$\sigma^2=1$ cup$^2$.
			\[
			\begin{array}{c|cccc}
				x & 1 & 2 & 3 & 4\\ \hline
				f & 2 & 5 & 4 & 1
			\end{array}
			\]
			
			\item
			\textbf{Mobile-data use (value--frequency table).}
			Daily mobile-data use, in hundreds of megabytes, was recorded for ten
			students. The theoretical values are $\mu=52$, $\sigma^2=4$, and $\sigma=2$.
			\[
			\begin{array}{c|cccc}
				x & 48 & 50 & 52 & 54\\ \hline
				f & 1 & 3 & 4 & 2
			\end{array}
			\]
			
			\item
			\textbf{Commute time (frequency buckets).}
			A random sample of commute times is grouped into the following intervals.
			Use class midpoints. The city model gives $\mu=30$ minutes and
			$\sigma=10$ minutes.
			\[
			\begin{array}{c|cccc}
				\text{Time (min)} & [10,20) & [20,30) & [30,40) & [40,50)\\ \hline
				f & 2 & 5 & 4 & 1
			\end{array}
			\]
			
			\item
			\textbf{Parcel mass (frequency buckets).}
			The masses of ten parcels are grouped below. Use class midpoints. Long-term
			records give $\mu=4$ kg and $\sigma^2=3$ kg$^2$.
			\[
			\begin{array}{c|cccc}
				\text{Mass (kg)} & [0,2) & [2,4) & [4,6) & [6,8)\\ \hline
				f & 1 & 4 & 4 & 1
			\end{array}
			\]
			
			\item
			\textbf{Package preparation time (mixed data).}
			A random sample of package preparation times, in minutes, was recorded in
			two different ways. Three observations were recorded individually:
			\[
			8,\quad 10,\quad 12.
			\]
			The remaining observations were summarized in a value--frequency table:
			\[
			\begin{array}{c|ccc}
				x & 9 & 10 & 11\\ \hline
				f & 2 & 3 & 2
			\end{array}
			\]
			Treat all the observations as one sample. Historical records give
			$\mu=10$ minutes and $\sigma^2=2$ minutes$^2$.
			
			\item
			\textbf{Customer service time (mixed data).}
			A random sample of customer service times, in minutes, was recorded using
			two different formats. Four observations were recorded individually:
			\[
			22,\quad 28,\quad 32,\quad 38.
			\]
			The remaining observations were grouped as follows:
			\[
			\begin{array}{c|cccc}
				\text{Time (min)} & [10,20) & [20,30) & [30,40) & [40,50)\\ \hline
				f & 1 & 2 & 2 & 1
			\end{array}
			\]
			Treat all the observations as one sample and use class midpoints for the
			grouped observations. The theoretical values are $\mu=30$ minutes and
			$\sigma=9$ minutes.
			
		\end{enumerate}
	\end{multicols}
	
	% ============================================================
	% SOLUTIONS
	% ============================================================
	
	\newpage
	\restoregeometry
	\vspace*{-15mm}
	
	\section{Solutions}
	
	\setlength{\columnseprule}{0.4pt}
	\begin{multicols}{2}
		
		% ============================================================
		% PROBLEM 1
		% ============================================================
		
		\subsection*{C7.1 -- Battery operating time}\phantomsection
		\label{subsec:c7-solution-batteries}
		
		\[
		n=5,\qquad
		\bar{x}=\frac{8+9+10+11+12}{5}=10.
		\]
		The sum of squared deviations is
		\[
		\sum(x-\bar{x})^2
		=4+1+0+1+4=10.
		\]
		Thus,
		\[
		s^2=\frac{10}{5-1}=2.5,\qquad
		s=\sqrt{2.5}\approx1.581.
		\]
		Since $\sigma=\sqrt{4}=2$,
		\[
		\boxed{\bar{x}=10,\quad s^2=2.5,\quad s\approx1.581.}
		\]
		Here $\bar{x}=\mu$, while $s^2<\sigma^2$ and $s<\sigma$; this sample has
		less spread than the population model.
		
		% ============================================================
		% PROBLEM 2
		% ============================================================
		
		\subsection*{C7.2 -- Local delivery time}\phantomsection
		\label{subsec:c7-solution-deliveries}
		
		\[
		n=6,\qquad
		\bar{x}=\frac{18+20+21+19+22+20}{6}=20.
		\]
		\[
		\sum(x-\bar{x})^2=4+0+1+1+4+0=10,
		\]
		so
		\[
		s^2=\frac{10}{6-1}=2,\qquad
		s=\sqrt{2}\approx1.414.
		\]
		The theoretical variance is $\sigma^2=1.5^2=2.25$. Therefore,
		\[
		\boxed{\bar{x}=20,\quad s^2=2,\quad s\approx1.414.}
		\]
		The sample mean equals $\mu$; both measures of sample spread are slightly
		below their theoretical values $2.25$ and $1.5$.
		
		% ============================================================
		% PROBLEM 3
		% ============================================================
		
		\subsection*{C7.3 -- Reusable cups}\phantomsection
		\label{subsec:c7-solution-cups}
		
		From the value--frequency table,
		\[
		\begin{aligned}
			n&=\sum f=2+5+4+1=12,\\
			\sum fx&=2(1)+5(2)+4(3)+1(4)=28,\\
			\sum fx^2&=2(1^2)+5(2^2)+4(3^2)+1(4^2)=74.
		\end{aligned}
		\]
		Hence
		\[
		\bar{x}=\frac{28}{12}=\frac73\approx2.333,
		\]
		\[
		\begin{aligned}
			s^2&=\frac{74-12(7/3)^2}{11}
			=\frac{26}{33}\approx0.788,\\
			s&=\sqrt{\frac{26}{33}}\approx0.888.
		\end{aligned}
		\]
		Since $\sigma=1$,
		\[
		\boxed{n=12,\quad\bar{x}\approx2.333,\quad s^2\approx0.788,
			\quad s\approx0.888.}
		\]
		All three sample statistics are below the corresponding theoretical values
		$2.5$, $1$, and $1$.
		
		% ============================================================
		% PROBLEM 4
		% ============================================================
		
		\subsection*{C7.4 -- Mobile-data use}\phantomsection
		\label{subsec:c7-solution-mobile-data}
		
		\[
		\begin{aligned}
			n&=1+3+4+2=10,\\
			\sum fx&=1(48)+3(50)+4(52)+2(54)=514,\\
			\sum fx^2&=1(48^2)+3(50^2)+4(52^2)+2(54^2)\\
			&=26452.
		\end{aligned}
		\]
		Therefore,
		\[
		\bar{x}=\frac{514}{10}=51.4,
		\]
		\[
		\begin{aligned}
			s^2&=\frac{26452-10(51.4)^2}{9}=3.6,\\
			s&=\sqrt{3.6}\approx1.897.
		\end{aligned}
		\]
		Thus
		\[
		\boxed{n=10,\quad\bar{x}=51.4,\quad s^2=3.6,\quad s\approx1.897.}
		\]
		The mean is $0.6$ below $\mu=52$; the variance and standard deviation are
		also slightly below $4$ and $2$.
		
		% ============================================================
		% PROBLEM 5
		% ============================================================
		
		\subsection*{C7.5 -- Commute time}\phantomsection
		\label{subsec:c7-solution-commute}
		
		The class midpoints of
		$[10,20)$, $[20,30)$, $[30,40)$, and $[40,50)$ are
		$15,25,35,45$. Thus,
		\[
		\begin{aligned}
			n&=2+5+4+1=12,\\
			\sum fx&=2(15)+5(25)+4(35)+1(45)=340,\\
			\sum fx^2&=2(15^2)+5(25^2)+4(35^2)+1(45^2)\\
			&=10500.
		\end{aligned}
		\]
		Using the midpoints,
		\[
		\bar{x}=\frac{340}{12}\approx28.333,
		\]
		\[
		\begin{aligned}
			s^2&=\frac{10500-12(340/12)^2}{11}
			=\frac{2600}{33}\approx78.788,\\
			s&\approx\sqrt{78.788}\approx8.876.
		\end{aligned}
		\]
		Therefore,
		\[
		\boxed{n=12,\quad\bar{x}\approx28.333,\quad s^2\approx78.788,
			\quad s\approx8.876.}
		\]
		The theoretical variance is $\sigma^2=10^2=100$. The grouped-sample
		estimates are below $\mu=30$, $\sigma^2=100$, and $\sigma=10$.
		
		% ============================================================
		% PROBLEM 6
		% ============================================================
		
		\subsection*{C7.6 -- Parcel mass}\phantomsection
		\label{subsec:c7-solution-parcels}
		
		The class midpoints of
		$[0,2)$, $[2,4)$, $[4,6)$, and $[6,8)$ are
		$1,3,5,7$. Therefore,
		\[
		\begin{aligned}
			n&=1+4+4+1=10,\\
			\sum fx&=1(1)+4(3)+4(5)+1(7)=40,\\
			\sum fx^2&=1(1^2)+4(3^2)+4(5^2)+1(7^2)=186.
		\end{aligned}
		\]
		Thus,
		\[
		\bar{x}=\frac{40}{10}=4,
		\]
		\[
		\begin{aligned}
			s^2&=\frac{186-10(4)^2}{9}=\frac{26}{9}\approx2.889,\\
			s&=\sqrt{\frac{26}{9}}\approx1.700.
		\end{aligned}
		\]
		Since $\sigma=\sqrt3\approx1.732$,
		\[
		\boxed{n=10,\quad\bar{x}=4,\quad s^2\approx2.889,
			\quad s\approx1.700.}
		\]
		The estimated mean equals $\mu$; the estimated variance and standard
		deviation are close to, but below, their theoretical values.
		
		% ============================================================
		% PROBLEM 7
		% ============================================================
		
		\subsection*{C7.7 -- Package preparation time}\phantomsection
		\label{subsec:c7-solution-package-preparation}
		
		There are three individual observations and seven observations in the
		value--frequency table. Therefore,
		\[
		n=3+(2+3+2)=10.
		\]
		
		Combining both representations,
		\[
		\begin{aligned}
			\sum x
			&=(8+10+12)+2(9)+3(10)+2(11)\\
			&=100.
		\end{aligned}
		\]
		Thus,
		\[
		\bar{x}=\frac{100}{10}=10.
		\]
		
		For the sum of squares,
		\[
		\begin{aligned}
			\sum x^2
			&=(8^2+10^2+12^2)
			+2(9^2)+3(10^2)+2(11^2)\\
			&=1012.
		\end{aligned}
		\]
		
		Therefore,
		\[
		\begin{aligned}
			s^2
			&=\frac{1012-10(10)^2}{10-1}\\
			&=\frac{12}{9}
			\approx1.333,
		\end{aligned}
		\]
		and
		\[
		s=\sqrt{\frac43}\approx1.155.
		\]
		
		Hence,
		\[
		\boxed{
			n=10,\quad
			\bar{x}=10,\quad
			s^2\approx1.333,\quad
			s\approx1.155.
		}
		\]
		
		The sample mean equals the theoretical mean $\mu=10$. The sample variance
		and standard deviation are below the theoretical values
		$\sigma^2=2$ and $\sigma=\sqrt2\approx1.414$.
		
		% ============================================================
		% PROBLEM 8
		% ============================================================
		
		\subsection*{C7.8 -- Customer service time}\phantomsection
		\label{subsec:c7-solution-customer-service}
		
		The grouped intervals have midpoints
		\[
		15,\quad25,\quad35,\quad45.
		\]
		
		There are four individual observations and six grouped observations, so
		\[
		n=4+(1+2+2+1)=10.
		\]
		
		Combining the individual observations with the grouped-data estimates,
		\[
		\begin{aligned}
			\sum x
			&=(22+28+32+38)
			+1(15)+2(25)+2(35)+1(45)\\
			&=300.
		\end{aligned}
		\]
		Therefore,
		\[
		\bar{x}=\frac{300}{10}=30.
		\]
		
		Similarly,
		\[
		\begin{aligned}
			\sum x^2
			&=(22^2+28^2+32^2+38^2)\\
			&\quad+1(15^2)+2(25^2)+2(35^2)+1(45^2)\\
			&=9686.
		\end{aligned}
		\]
		
		Hence,
		\[
		\begin{aligned}
			s^2
			&=\frac{9686-10(30)^2}{10-1}\\
			&=\frac{686}{9}
			\approx76.222,
		\end{aligned}
		\]
		and
		\[
		s=\sqrt{\frac{686}{9}}\approx8.731.
		\]
		
		Therefore,
		\[
		\boxed{
			n=10,\quad
			\bar{x}=30,\quad
			s^2\approx76.222,\quad
			s\approx8.731.
		}
		\]
		
		The estimated sample mean equals $\mu=30$. The estimated standard deviation
		is slightly below the theoretical value $\sigma=9$. Because part of the
		sample is grouped, the statistics obtained using class midpoints are
		estimates.
		
	\end{multicols}
	
\end{document}